\section*{Chapter \ref{ch:qm:linear-models-under-stress} --- Linear Models under Stress}

\begin{solution}{exo:qm:linear-models-under-stress:1}
$R_j^2 = 0.81$, so $\mathrm{VIF} = 1/0.19 = 5.3$ and each \omterm{def:qm:estimation:estimator}{standard error} grows by $\sqrt{5.3} = 2.3$.
\end{solution}

\begin{solution}{exo:qm:linear-models-under-stress:2}
$\lambda = 1/(1 + 0.25) = 0.8$, so OLS estimates $0.8 \times 2 = 1.6$.
\end{solution}

\begin{solution}{exo:qm:linear-models-under-stress:3}
$\sum_ih_{ii} = \operatorname{tr}H = \operatorname{tr}\bigl(X(X^\top X)^{-1}X^\top\bigr) = \operatorname{tr}\bigl((X^\top X)^{-1}X^\top X\bigr) = \operatorname{tr}I_k = k$.
\end{solution}

\begin{solution}{exo:qm:linear-models-under-stress:4}
With $X^\top X = nI$ the objective separates: $\frac1{2n}\lvert y - X\beta\rvert^2 = \mathrm{const} + \frac12\sum_j(\beta_j - b_j)^2$ with $b = X^\top y/n$ the
least-squares coefficient. Each coordinate minimises $\frac12(\beta_j - b_j)^2 + \lambda|\beta_j|$, whose solution is $\operatorname{sign}(b_j)\max(|b_j| -
\lambda, 0)$: soft thresholding.
\end{solution}

\begin{solution}{exo:qm:linear-models-under-stress:5}
The level-error variance is $-\gamma(1) = 0.0196$ (standard deviation 0.14), so $\hat\lambda = (0.199 - 2 \times 0.0196)/0.199 = 0.80$. Weekly, the true
variance is $5 \times (0.199 - 0.0392) = 0.80$ against the same $0.0392$ of error: $\lambda_5 = 0.80/0.839 = 0.95$.
\end{solution}

\begin{solution}{exo:qm:linear-models-under-stress:6}
With $X^\top X \approx n\begin{psmallmatrix}1 & \rho\\ \rho & 1\end{psmallmatrix}$, $\Var(\hat\beta_1) = \sigma^2/(n(1 - \rho^2)) = 16.9\sigma^2/n$ while
$\Var(\hat\beta_1 + \hat\beta_2) = 2\sigma^2/(n(1 + \rho)) = 1.02\sigma^2/n$. The sum, the exposure to the common direction, is pinned down about four times
more precisely (in standard deviation) than either beta; the difference, with variance $2\sigma^2/(n(1 - \rho)) = 66.7\sigma^2/n$, is where the
noise lives.
\end{solution}

\begin{solution}{exo:qm:linear-models-under-stress:7}
\texttt{panel\_truth()}: true standard deviation 0.105; classical 0.026 and White 0.026; clustered by firm 0.101; Fama--MacBeth 0.023 (0.022 with three
Newey--West lags). The firm effect is common to every month's slope, so neither the classical formula nor the time-series spread of the monthly
slopes sees it; clustering by firm does.
\end{solution}

\begin{solution}{exo:qm:linear-models-under-stress:8}
Random folds put days from the future into the training set of every test day, and with twenty-day targets overlapping by nineteen days,
neighbouring observations share most of their target: the model is tested on returns it has effectively seen. Use expanding (or rolling) folds in
time order with a purge of at least twenty days between training and test, and expect a far lower $R^2$; also count the penalties and
specifications tried (chapter~12).
\end{solution}

\begin{solution}{pb:qm:linear-models-under-stress:1}
\textbf{1.} 0.68.
\textbf{2.} $\lambda = 0.16/(0.16 + 2 \times 0.14^2) = 0.80$, so on average $0.80 \times 0.85 = 0.683$.
\textbf{3.} Mean 0.683, standard deviation 0.017.
\textbf{4.} The noise in the regressor biases the slope; more data shrink the variance around the wrong value, not the bias.
\textbf{5.} 0.086 points a day against 0.049 at the true ratio, 74\% more.
\textbf{6.} $-0.098$; $\hat\lambda = 1 + 2 \times (-0.098) = 0.80$.
\textbf{7.} $0.68/0.80 = 0.85$, residual 0.049.
\textbf{8.} $\lambda_5 = 0.95$; the weekly ratio is 0.80, residual 0.053.
\textbf{9.} Corrected: mean 0.862, standard deviation 0.095. Weekly: mean 0.810, standard deviation 0.022.
\textbf{10.} 0.74; it assumes the error variances in the bond and the future are equal, while here the future's is 0.039 and the bond's residual
0.0025.
\textbf{11.} $0.16/(0.16 + 2 \times 0.07^2) = 0.94$.
\textbf{12.} It removes the problem: no level error, no attenuation, and the daily OLS ratio is unbiased.
\textbf{13.} Monthly, $\lambda_{21} = 21 \times 0.16/(21 \times 0.16 + 0.039) = 0.99$, almost no bias, but only 24 observations in two years, so a noisy
ratio; the weekly regression is the compromise.
\textbf{14.} No: collinearity is a lack of independent variation, not noise in a regressor. Aggregating does not separate the factors; only reporting
their sum, penalising, or finding data where they move apart does.
\textbf{15.} Compare with the duration-based ratio, test the hedged P\&L out of sample, and check the ratio's stability across periods and
horizons.
\textbf{16.} The corrected or duration-based ratio near 0.85, cross-checked by the weekly 0.80; not the daily OLS 0.68.
\textbf{17.} Record synchronous mids for the future at the bond's marking time.
\textbf{18.} When the goal is prediction from the noisy regressor itself (the attenuated slope is then the best linear predictor), rather than the
structural relation.
\textbf{19.} Named result: \emph{the hedge ratio that shrank}: stale, bid-ask-noisy future prices attenuate the daily OLS hedge ratio by 0.80, from
0.85 to 0.68; the ratio corrected from the future's own first autocorrelation is 0.85; the residual risk is 0.086 points a day at 0.68 against 0.049
at 0.85 (0.053 at the weekly 0.80).
\textbf{20.} It biases the slope toward zero, while noise in the dependent variable only adds variance.
\end{solution}

\begin{solution}{iq:qm:linear-models-under-stress:1}
Each coefficient's variance is inflated by $1/(1 - 0.97^2) = 17$; the individual coefficients swing and can flip sign, while their sum is well
determined. Report the sum or the principal-component exposure, penalise (ridge), or drop one factor.
\iqlookfor{The VIF, the stable combination, and a remedy chosen for the question asked.}
\end{solution}

\begin{solution}{iq:qm:linear-models-under-stress:2}
Attenuation from noise in the regressor: asynchronous or stale prices, bid-ask bounce in the future's prints. Also a changing relation (the
sample spans different regimes), or the wrong horizon. Check the future's first autocorrelation, rerun on weekly changes, and use synchronous
prices.
\iqlookfor{\omterm{def:qm:linear-models-under-stress:eiv}{Errors-in-variables} as the first suspect, with a diagnostic and a fix.}
\end{solution}

\begin{solution}{iq:qm:linear-models-under-stress:3}
Ridge shrinks along the small singular directions and spreads weight across a correlated block; the \omterm{def:qm:linear-models-under-stress:penalties}{lasso} sets coefficients to zero and tends to
keep one member of a block, somewhat arbitrarily; the \omterm{def:qm:linear-models-under-stress:penalties}{elastic net} keeps groups together while still selecting.
\iqlookfor{Shrinkage versus selection, and the behaviour within correlated groups.}
\end{solution}

\begin{solution}{iq:qm:linear-models-under-stress:4}
Folds in time order, training before testing, with a purge gap of at least the target's horizon (twenty days) so overlapping targets do not leak,
and possibly an embargo after the test block; standardise and choose features inside each training fold only.
\iqlookfor{Time order, purge for overlap, and no preprocessing on the full sample.}
\end{solution}

\begin{solution}{iq:qm:linear-models-under-stress:5}
When the regressors and residuals have persistent firm effects: every period's slope shares the same error, so the time-series standard deviation
understates the uncertainty. Use \omterm{def:qm:estimation:estimator}{standard errors} clustered by firm (or by firm and period when both effects exist).
\iqlookfor{Which dependence each method handles: Fama--MacBeth handles time effects, clustering by firm handles firm effects.}
\end{solution}

\begin{solution}{iq:qm:linear-models-under-stress:6}
The coefficient on $X_2$ in the regression on $(X_1, X_2)$ equals the slope of $M_1y$ on $M_1X_2$, the parts orthogonal to $X_1$; proof by applying
$M_1$ to the normal equations. On a desk: a stock's beta to a sector after removing the market is the slope of market-residual returns; a signal's
marginal value after existing signals is its fit on residualised returns.
\iqlookfor{The projection argument and a concrete residualisation.}
\end{solution}
